%! Direct, Inverse \& Joint Variation — Unit 4, Lesson 4.7 — Exit Ticket (Answer Key)
\documentclass[10pt]{article}
\usepackage{algebra2-article}
\usepackage{algebra2-key}   % pulls in -boxes; do NOT also load -boxes
\newcommand{\TallMath}[1]{$\displaystyle #1\rule[-1.4em]{0pt}{3.2em}$}

\begin{document}

\pageheader{Unit 4, Lesson 4.7}{Exit Ticket --- Answer Key}
\namedateperiod

\vspace{0.06in}

No notes. Run \emph{both} tests before you name anything.

\vspace{0.04in}

\begin{enumerate}[leftmargin=1.8em, itemsep=9pt]
  \item \textbf{Name it, then use it.}
    \begin{center}
    \renewcommand{\arraystretch}{1.25}
    \begin{tabular}{|l|c|c|c|c|}
    \hline
    \rowcolor{forestbg}
    $x$ & $2$ & $3$ & $4$ & $6$ \\ \hline
    $y$ & $18$ & $12$ & $9$ & $6$ \\ \hline
    $y\div x$ & \ans{$9$} & \ans{$4$} & \ans{$2.25$} & \ans{$1$} \\ \hline
    $x\cdot y$ & \ans{$36$} & \ans{$36$} & \ans{$36$} & \ans{$36$} \\ \hline
    \end{tabular}
    \end{center}
    \vspace{-2pt}
    \begin{enumerate}[label=(\alph*), leftmargin=2.1em, itemsep=4pt]
      \item This is an \ans{inverse} variation, because the \ans{product} row is constant.
      \item $k=$ \ans{$36$} \quad Equation: $y=$ \ans{$\dfrac{36}{x}$}
      \item Find $y$ when $x=8$: \ans{$4.5$} \quad Find $x$ when $y=3$: \ans{$12$}
      \item The graph would have a vertical asymptote at \ans{$x=0$} and a horizontal asymptote
            at \ans{$y=0$}.
    \end{enumerate}

  \item \textbf{Correct the reasoning.} A classmate looks at the table below and says: \emph{``$y$ goes
        down every time $x$ goes up, so $y$ varies inversely as $x$.''}
    \begin{center}
    \renewcommand{\arraystretch}{1.25}
    \begin{tabular}{|l|c|c|c|c|}
    \hline
    \rowcolor{forestbg}
    $x$ & $1$ & $2$ & $3$ & $4$ \\ \hline
    $y$ & $12$ & $9$ & $6$ & $3$ \\ \hline
    $x\cdot y$ & \ans{$12$} & \ans{$18$} & \ans{$18$} & \ans{$12$} \\ \hline
    \end{tabular}
    \end{center}
    \vspace{-2pt}
    The correct verdict is \ans{neither}. Write one or two sentences telling your classmate what test
    they skipped and what it showed. \ansline{``$y$ goes down'' is not the test --- the test is whether
    the product $xy$ is the same for \emph{every} pair.}
    \ansline{Here the products are $12, 18, 18, 12$: two of them agree, but not all four, so this is
    neither variation. (It is the line $y=15-3x$.)}

  \item \textbf{SOL-style multiple choice.} The time $t$ needed to fill a pool varies inversely as the
        rate $r$ of the pump. A pump running at $8$ gallons per minute fills the pool in $45$ minutes.
        How long would a $12$ gallon-per-minute pump take?
        \vspace{-6pt}
        \begin{enumerate}[label=\Alph*., leftmargin=1.9em, itemsep=1pt]
          \item $30$ minutes \quad \textcolor{keyred}{\textbf{$\leftarrow$ correct}}
          \item $67.5$ minutes
          \item $3.75$ minutes
          \item $22.5$ minutes
        \end{enumerate}
\end{enumerate}

\begin{teachernote}
\textbf{Item 1.} The quotient row is there on purpose --- a student who fills in only the product row
guessed rather than tested. Part (d) closes the unit: $y=\frac{36}{x}$ is the 4.4 parent stretched by
$36$, so the asymptotes are the axes and $k$ cannot move them.

\vspace{2pt}
\textbf{Item 2 is the graded-language item.} The products $12, 18, 18, 12$ are built so two of them
match, catching students who check a row or two and stop. Grade the \emph{reason}: full credit names
the constant-product test and reports that it fails across all four pairs. ``It's not inverse'' with no
test cited earns nothing.

\vspace{2pt}
\textbf{Item 3.} $k=rt=8\cdot45=360$ gallons, so $t=\frac{360}{12}=30$ --- answer \textbf{A}. Each
distractor is a named error: \textbf{B} sets up a direct proportion; \textbf{C} divides $45$ by the new
rate without finding $k$; \textbf{D} halves $45$ because the rate ``went up.'' Sense check: a faster
pump must take \emph{less} time, which kills B on sight.

\vspace{2pt}
\textbf{Sorting the pile:} got it / named it from one row / used the wrong test / computed correctly but
could not say what $k$ means. Catch the last group before the unit test, where interpretation carries
as many points as the arithmetic.
\end{teachernote}

\end{document}
